\paragrafos{Παραπληρωματικές, συμπληρωματικές και κατακορυφήν γωνίες}{Β.1.8, «Παραπληρωματικές και Συμπληρωματικές γωνίες – Κατακορυφήν γωνίες»}

\begin{anagke}
Δύο μονοπάτια του κήπου διασταυρώνονται σαν ένα μεγάλο «Χ». Η Αγγελική μετράει τη μία γωνία και βρίσκει $48^\circ$. «Πόσες είναι οι άλλες τρεις;» ρωτάει. Ο Νίκος μετράει όλες: $48^\circ$, $132^\circ$, $48^\circ$, $132^\circ$. «Οι απέναντι είναι ίσες! Και οι διπλανές κάνουν μαζί $180^\circ$». «Είναι τυχαίο;» ρωτάει η Πολυξένη. «Μήπως αν κάνω το μονοπάτι πιο λοξό, θα αλλάξει;». Δοκιμάζουν με άλλη γωνία, $70^\circ$, και βρίσκουν $70^\circ$, $110^\circ$, $70^\circ$, $110^\circ$.

Η Ρήνη ρωτάει για μια γωνία που τη συμπληρώνουμε μέχρι την ορθή: ένα παρτέρι σχήματος γωνίας $35^\circ$ και το υπόλοιπο μέχρι τη γωνία του φράχτη των $90^\circ$ θα φυτευτεί με τριαντάφυλλα: «Πόσες μοίρες μένουν;». Ο Στέλιος υπολογίζει: «Εδώ το άθροισμα είναι $90^\circ$ και όχι $180^\circ$».

Έτσι η παρέα συναντά τρεις έννοιες: τις παραπληρωματικές γωνίες, που αθροίζουν $180^\circ$, τις συμπληρωματικές, που αθροίζουν $90^\circ$, και τις κατακορυφήν, που βρίσκονται απέναντι όταν δύο ευθείες τέμνονται. Θα βρουν και γιατί οι κατακορυφήν γωνίες είναι ίσες, χωρίς να χρειάζονται μοιρογνωμόνιο.
\end{anagke}

\begin{lexi}{παραπληρωματική, συμπληρωματική, κατακορυφήν}{παραπληρωματική, συμπληρωματική, κατακορυφήν}
Η \lexiword{συμπληρωματική} γωνία «συμπληρώνει» μια γωνία μέχρι την ορθή. Η \lexiword{παραπληρωματική} «παρα-πληρώνει», δηλαδή «συμπληρώνει ως τη μισή στροφή»: ως τις $180^\circ$. Οι \lexiword{κατακορυφήν} γωνίες είναι αυτές που βρίσκονται «κατά κορυφήν», δηλαδή απέναντι, με κοινή κορυφή.
\end{lexi}

\begin{metafrasi}
Δύο γωνίες που αθροίζουν $90^\circ$. & συμπληρωματικές: $\omega+\varphi=90^\circ$\\
Δύο γωνίες που αθροίζουν $180^\circ$. & παραπληρωματικές: $\omega+\varphi=180^\circ$\\
Οι απέναντι γωνίες δύο τεμνόμενων ευθειών. & κατακορυφήν: ίσες\\
Η συμπληρωματική της $35^\circ$. & $90^\circ-35^\circ=55^\circ$\\
Η παραπληρωματική της $130^\circ$. & $180^\circ-130^\circ=50^\circ$\\
\end{metafrasi}

\tmima{Η θεωρία}

\begin{orismos}{συμπληρωματικές, παραπληρωματικές}
Δύο γωνίες λέγονται \textbf{συμπληρωματικές} αν το άθροισμά τους είναι $90^\circ$ και \textbf{παραπληρωματικές} αν το άθροισμά τους είναι $180^\circ$. Η καθεμία λέγεται συμπλήρωμα (ή παραπλήρωμα) της άλλης.
\end{orismos}

Για παράδειγμα η συμπληρωματική της $35^\circ$ είναι η $55^\circ$ και η παραπληρωματική της $130^\circ$ η $50^\circ$. Δύο εφεξής παραπληρωματικές γωνίες έχουν τις μη κοινές πλευρές τους αντικείμενες ημιευθείες. Δύο συμπληρωματικές γωνίες είναι πάντα οξείες, ενώ από δύο παραπληρωματικές η μία είναι αμβλεία και η άλλη οξεία (ή και οι δύο ορθές).

\tmimasmall{Κατακορυφήν γωνίες}

Όταν δύο ευθείες τέμνονται σχηματίζονται τέσσερις γωνίες. Δύο από αυτές που δεν είναι διπλανές λέγονται \textbf{κατακορυφήν}.

\begin{center}
\begin{tikzpicture}[scale=1.1]
\draw[thick,{Stealth}-{Stealth}] (-2.5,-1)--(2.5,1); \draw[thick,{Stealth}-{Stealth}] (-1.2,-1.8)--(1.2,1.8);
\fill (0,0) circle (2pt); \node[below right] at (0.05,0) {$\mathrm{Ο}$};
\node[gblue] at (1.1,0.18) {$\omega$}; \node[gblue] at (-1.1,-0.18) {$\omega$}; \node at (0.18,0.9) {$\varphi$}; \node at (-0.18,-0.9) {$\varphi$};
\end{tikzpicture}
\end{center}

\begin{idiotita}{κατακορυφήν γωνίες}
Οι κατακορυφήν γωνίες είναι \textbf{ίσες}.
\end{idiotita}

Γιατί: η $\omega$ και η διπλανή της $\varphi$ είναι παραπληρωματικές ($\omega+\varphi=180^\circ$). Η κατακορυφήν της $\omega$, έστω $\omega'$, είναι κι αυτή παραπληρωματική της $\varphi$ ($\omega'+\varphi=180^\circ$). Άρα $\omega=180^\circ-\varphi=\omega'$.

\begin{prosoxi}
\begin{itemize}[leftmargin=1.2em,itemsep=0.3ex,topsep=2pt]
\item Συμπληρωματικές είναι $90^\circ$, παραπληρωματικές $180^\circ$. Μην τις μπερδεύεις: η «παραπληρωματική» είναι η μεγαλύτερη λέξη και το μεγαλύτερο άθροισμα ($180^\circ$).
\item Δεν υπάρχει συμπληρωματική γωνία για γωνία μεγαλύτερη ή ίση με $90^\circ$ (θα έβγαινε μηδενική ή αρνητική).
\item Οι διπλανές γωνίες δύο τεμνόμενων ευθειών δεν είναι ίσες (εκτός αν είναι ορθές): είναι παραπληρωματικές.
\item Δύο γωνίες που αθροίζουν $180^\circ$ δεν είναι αναγκαστικά εφεξής: αρκεί το μέτρο.
\end{itemize}
\end{prosoxi}

\tmima{Λυμένα παραδείγματα}

\paradeigma{Συμπλήρωμα και παραπλήρωμα}
\tixero Η γωνία $35^\circ$ και η γωνία $130^\circ$.\par
\tithelo Η συμπληρωματική της πρώτης και η παραπληρωματική της δεύτερης.\par
\possindeo $90^\circ-35^\circ=55^\circ$. $180^\circ-130^\circ=50^\circ$.\par
\apantisi $55^\circ$ και $50^\circ$.\par
\elegxos $35+55=90$ και $130+50=180$.

\paradeigma{Δύο τεμνόμενες ευθείες}
\tixero Δύο ευθείες τέμνονται και μία γωνία είναι $48^\circ$.\par
\tithelo Οι άλλες τρεις γωνίες.\par
\possindeo Η κατακορυφήν της είναι $48^\circ$. Οι διπλανές είναι παραπληρωματικές της: $180^\circ-48^\circ=132^\circ$, και η απέναντί τους επίσης $132^\circ$.\par
\apantisi $48^\circ$, $132^\circ$, $132^\circ$.\par
\elegxos Το άθροισμα και των τεσσάρων: $48+132+48+132=360$.

\paradeigma{Παραπληρωματικές με λόγο}
\tixero Δύο παραπληρωματικές γωνίες έχουν λόγο $2:3$.\par
\tithelo Να τις βρω.\par
\possindeo Αν η μία είναι $2x$ και η άλλη $3x$, $2x+3x=180$, $5x=180$, $x=36$. Οι γωνίες είναι $72^\circ$ και $108^\circ$.\par
\apantisi $72^\circ$ και $108^\circ$.\par
\elegxos $72+108=180$ και $\frac{72}{108}=\frac{2}{3}$.

\paradeigma{Εξίσωση με κατακορυφήν}
\tixero Οι κατακορυφήν γωνίες είναι $2x+10$ και $3x-20$ μοίρες.\par
\tithelo Το $x$ και οι γωνίες.\par
\possindeo Ίσες: $2x+10=3x-20$, άρα $x=30$ και οι γωνίες είναι $2\cdot30+10=70^\circ$.\par
\apantisi $x=30$ και γωνίες $70^\circ$.\par
\elegxos $3\cdot30-20=70$.

\tmima{Ασκήσεις}

\zoni{Το κατάλαβα}{συμπληρωματικές, παραπληρωματικές, κατακορυφήν}

\ex Βρες το συμπλήρωμα των γωνιών: $35^\circ$, $62^\circ$, $8^\circ$, $44^\circ30'$.

\ex Βρες το παραπλήρωμα των γωνιών: $130^\circ$, $90^\circ$, $5^\circ$, $73^\circ15'$.

\ex Δύο ευθείες τέμνονται και μία γωνία είναι $48^\circ$. Βρες τις άλλες τρεις.

\ex Δύο ευθείες τέμνονται και μία γωνία είναι ορθή. Βρες τις άλλες.

\ex Βρες γωνία που είναι ίση με το συμπλήρωμά της. Βρες γωνία ίση με το παραπλήρωμά της.

\ex Δύο παραπληρωματικές γωνίες έχουν λόγο $2:3$. Βρες τες.

\ex Σωστό ή λάθος; α) Οι κατακορυφήν γωνίες είναι ίσες. \quad β) Οι διπλανές γωνίες δύο τεμνόμενων ευθειών είναι συμπληρωματικές. \quad γ) Υπάρχει συμπληρωματική της $100^\circ$. \quad δ) Δύο συμπληρωματικές γωνίες είναι οξείες. \quad ε) Οι παραπληρωματικές γωνίες έχουν άθροισμα $180^\circ$.

\zoni{Το χρησιμοποιώ}{διασταυρώσεις, παρτέρια, καθρέφτες}

\ex \textbf{Η διασταύρωση.} Δύο μονοπάτια του κήπου τέμνονται και μια γωνία είναι $70^\circ$. Πόσες μοίρες έχουν οι άλλες γωνίες; Πόσες διαφορετικές τιμές υπάρχουν;

\ex \textbf{Τα τριαντάφυλλα.} Στη γωνία $90^\circ$ του φράχτη φυτεύεται ένα παρτέρι με γωνία $35^\circ$. Πόση γωνία μένει για τα τριαντάφυλλα;

\ex \textbf{Ο καθρέφτης.} Ένα χέρι του κήπου σχηματίζει γωνία $\omega=2x+10$ και η κατακορυφήν της είναι $3x-20$. Βρες το $x$ και την $\omega$.

\ex \textbf{Η σκάλα.} Μια σκάλα ακουμπά σε έναν τοίχο σχηματίζοντας γωνία $72^\circ$ με το έδαφος. Πόση γωνία σχηματίζει με τον τοίχο; (Ο τοίχος είναι κάθετος στο έδαφος.)

\ex \textbf{Το ψαλίδι.} Ένα ψαλίδι ανοίγει ώστε οι λεπίδες να σχηματίζουν γωνία $40^\circ$. Πόση είναι η γωνία των λαβών (κατακορυφήν);

\zoni{Το σκέφτομαι}{γιατί ίσες}

\ex Εξήγησε γιατί οι κατακορυφήν γωνίες είναι ίσες, χρησιμοποιώντας τα παραπληρώματα.

\ex Δύο γωνίες είναι ίσες. Είναι απαραίτητα κατακορυφήν; Δώσε αντιπαράδειγμα.

\ex Αν η $\omega$ είναι οξεία, τι είδους είναι η παραπληρωματική της; Και η συμπληρωματική της;

\ex Οι διχοτόμοι δύο εφεξής παραπληρωματικών γωνιών: πόση γωνία σχηματίζουν μεταξύ τους; (δοκίμασε με $\omega=50^\circ$).

\zoni{Για το σπίτι}{ψάξε Χ}

\ex Σχεδίασε δύο τεμνόμενες ευθείες, μέτρησε και τις τέσσερις γωνίες και έλεγξε ότι οι κατακορυφήν είναι ίσες.

\ex Βρες στο σπίτι ένα ζευγάρι τεμνόμενων ευθειών (π.χ. ψαλίδι, παραθύρων) και μέτρησε τις γωνίες.
